\documentclass[../main.tex]{subfiles}

\graphicspath{{\subfix{../images/}}}



\begin{document}



\clearpage
\thispagestyle{empty}

\vspace*{7cm}
\section{Power Supplies}
\vspace*{1.5cm}






\subsection{Introduction}

Power supplies are a fundamental component in electronic circuits and integrated circuits, as they provide the energy necessary for their operation.

Each type of electronic circuit requires specific power supply characteristics: voltage, current, stability, and noise level can vary significantly depending on the application. For this reason, there are various types of power supplies, designed to meet different needs.

For example, some circuits require a stable and precise DC voltage, while others can tolerate greater variations or require higher power. Consequently, the choice of power supply depends strictly on the type of circuit to be powered and the required performance. 

In the following sections we will be introducing and then analyzing the basics concepts and aspects of power supplies and moreover of power electronics in general.

We will begin with a general overview moving then to classical power supply introducing then voltage regulators to move then to switching power supplies.

\subsubsection{Some useful definitions}

\subsection{Specifications}

As we can imagine, given that we’re talking about the embedded field, various devices can be powered by a wide variety of energy sources, such as wall sockets, solar panels, batteries, and dynamos—each of which can provide energy in different forms, which is why it must be converted, so to speak, to be compatible with the system we want to power.  

Starting by looking at the specifications, we must think of a power supply as a circuit that transforms the energy source into electrical quantities compatible with the circuit to be powered and that must satisfy electrical and other specifications.

In particular, electronic circuits—which are usually powered by direct current—require a constant, predetermined voltage and draw current based on their usage. 

Let’s now delve into more specific details and analyze the required specifications, which include: input specifications, output specifications, power dissipation, and transient behavior.

Let’s begin by analyzing them one by one, starting with the input specifications—specifically, the voltage.

\subsubsection{Electrical specifications}


\subsubsection*{Input voltage}

Iniziamo ora con l'analisi di una delle variabili in ingresso ovvero il vlaore della tensione. Come detto in rpecedneza per questa tipolgia di sistemi, dalle specifiche dateci dal costruttore o definite da noi stessi, dovremo alimentare il circuit embedded ad una data tensione e lasciare che, in base al suo utilizzo, assorba la corrente necessaria al suo funzioanmento, motivo per cui la tensione é si un valroe importante ma soprattuto in sucora dal sistema di alimentazione non quanot in entrata.

Come possiamo immagianre i istemai di aliemntazione devono poiter funzioanre in tutto il mondo otivo per cui, nella stramaggioranza dei casi, sono disegnati per poter accettare in ingresso tutti i valori di tensioni che le reti di sitribuziini nazionali forniscosocno con anche un ampio magrine di errore per aggiungere un maggiore livello di sicurezza.

Per fare delgi esempi possiamo trovare tra i valori più comuni vlaori nominali di tensione in ingresso troviamo:

\begin{enumerate}
	\item[\textbf{-}] $V_{in}$ = 3.7 V $\pm$ 2\%
	\item[\textbf{-}] $V_{in}$ = 230 $\unit{V_{RMS}}$ $\pm$ 10\%
	\item[\textbf{-}] $V_{in}$ = 110 $\unit{V_{RMS}}$ 10\%
\end{enumerate}

Inoltre i dispistivi accettano in ingreso un determianto range di tensini co una così grnde toleranza in quanto la tensione in ingresso, che potremmo anche intepretare come un segnale, é soggetta a rumori e potrebbe variare molto rispetto al valore in cui si deve toravre. Per esempio la tensione di una batteria di una macchina potrebbe non essere sempre  uguale e variare in base all'utilizzo della stessa (power on fino a 60 V) in maniera stocastica.

Oltre alla tensione in ingresso dobbiamo, ripensando semrpe all'intup come ad un segnale,dobbiamo considerare alal frequenza a cui il segnale viene fornito.

Tendenzialmente questo valore, per gli apparechhi che si connettono alla rete pubblica ha un valore di 60 o 50 Hz pm 5 perc in base alla zone in cui ci si trova nel mondo.

Infine tra le spesiciche di inpiut vanno anche considerati i transienti, ovvero i comportamnti che assume il nostro unpu, anche seppuer per brevi periodi, al passaggio tra uno stato di regime ad un'altro. Ne possono essere un'sempio:

\begin{enumerate}
	\item[\textbf{-}] \textbf{Sag:} short reduction (few second) reduction of the nominal voltage
	\item[\textbf{-}] \textbf{Brownout:} long reduction (minutes or hours) reduction of the nominal voltage
	\item[\textbf{-}] \textbf{Micro-interruptions:} changes in the voltage value that are repeated in a short period of time.
\end{enumerate}

\subsubsection*{Output voltage}

Adesso che abbiamo vista le specifiche di tensione in ingresso possiamo andare a vedere il voltaggi odi uscirsca. In aprtcioarem partiamo dal presupposto che un singolo apparato di alientazione può generare in uscita una o più tensioni in base all'applicazione che deve eseguire.

A differenza di quanti visot per le tensini di ingresso i valori nominali e le tolleranze dei valori in uscita dal nostro apprato devono essere minori. Tra i vlaori possiamo trovare: VRMS V di picco ed il vlaore medio.

In partcioalre considerando ora la tensinone in usncita come uns segnale ci dovremo anche preiccpura di eventiale rumore introdotto oppure di effetti quali il ripple ovvero della ripetiizone periodica di un certo tipo di effetto non desiderato. 

Come appena detto, e come possiamo immagianre, le specifiche di outout deono essere maggiromane stringenti in quanit anche l;alimentazione é un fattoe chairo per il corretto funzioanet del nostro apparato tanto quanto il codice e il corretto hardware. 

Troviamo ora altre specifiche di output come la sua stabilità ed alla regolazione della linea in uscita che limita la pssobilità all tensione di variare nel tempo.
In partcioalre dobbiamo teenre sotto comtroll fenomnei quali la varaizione della tensione di usicta a seguoto della variazione di quella in ingresso.

Possiamo esprimere quetsa relazione come:

\begin{equation}
	\Delta V_{out} \big|_{\Delta_{V_{in}}}
\end{equation}

notiamo inoltre come, solitamente, la differenza di tensione che deve essere pplicata in ingresso deve essere significativa per produrre delle variazioni sulla tensione di usctia di comuqnue circa due ordini di grandezza inferiori nel caso in cui ci si torvi davanti ad un buon design.

Altri fattori che impattano sono la tmepereatura e la vita del nostro pparato.
In particoalre possiao risocntrare discotamenti nell'ordine dei ppm/grado centrigrado nella tensione di sufita del nostro alimentaore mentre si possoino vedere delle varaizioni nella tensione di uscita di circa mV/mese quando sui consideran la vita del nostro circuito.

Infine analizziamo anche gli aspetti dinamici del nosti circuit.

La maggior parte di tlai fenamone diopende dal carico applciato e da dinamiche particolai quali la tiplogia del carico e le tempistiche in cui ciò avviene.

In aprtcioalre essendo l'amiemntatore uns sistema a retroazione (feedback) i suoi poli determinano fattori quali velocità di risposta, stabilità, presenza di oscillazioni.
Quindi poli ben progettati portano a risposte veloci e stabili mentre progetti non corretti presentano fenomeni quali overshoot, oscillazioni, instabilità. Quesit fattori sono findamentlai per capire dopo quanti tempo, per esempio, il segnale di suita torna ad eseere accettaboel diopo una certo perioddo tanaisotrio in quato alcune tensini o comporamtnri non potrebbero essere ben accettati del dispositivoalimentato.

Altri fenomeni si possono verificare a seguoto di una cambiamento della corrente assorbita e fornita in sucita, più velovmente la tensione resta nel valore corretto meglio é per il nostro istema. in questo caso la corrente cambia bruscamente (step load) da 10 mA a 200 mA istantaneamente generando un calo (dip) o un picco (overshoot) della tensione.

% FOTO

In altri termini vanno anche considerati undershoot (minimo) overshoot (massimo) ovvero ivalori massimoi che la tensione può assumere durante una fanse trnaistoria, non si possono superare alcuni limiti dato che si potrebbe incappare in prble iqual iil reset del microcontrollore.

Altri paremetri che rendono ancora più strigneti le spcifiche di outpout possono essere varibili come laposisiblit`q dell;utente di modficare trmaite un nsleetotre o alktro dsispitivio (partitore resistivo) la tensione di uscita.
Questo aggiungi complicazioni ijn quanot tutto quello che abbiao detto sopra deve valere per tutto il range di tensioni disponibile in usita.

Infine citiamo latri parametri quali l'hold-up time ovvero il tempo per cui l'output resta valido dal momento in cui viene a mancare la tensione in ingresso. Questo parametr opuò essere utile, nel caos sià geande, per avvisare in tempo il microcontroller di salvare dati importanti prima di spegnersi dato che é avvenuto una qalche sorta di problema.

Questo vlaore sipende dall'utilizzo del carico e dalla tipologia del condesnatore in uscita per esempio se l'uscita di un alimentatore resta a 5 V per un tempo di 10 ms prima di scendere sotto la soglia indicata e andare a zero quello sarà il tempo di hold-on.

Dovessimo andare afare un piccolo recup un buon aliemnatore, il che non vuole dire che vada bene per qualsiaisi appliczone, deve essere:
, mantenere la tensione corretta (output voltage) reagire bene ai cambiamenti (dynamic behavior) recuperare velocemente (recovery time), evitare picchi pericolosi (overshoot/undershoot), essere regolabile (adjustable range), resistere a brevi blackout (hold-up time)


\subsubsection*{Output current}

terminiamo ora la nostro vista sulle pscifiche di output andando a venedere in ultima istanza la corrente in uscita e le sue specifiche / limitaizoni.

Prima di tutto dibbiamo dire che la corrente di uscita, ovvero qualle assorbiata dal nostro sistema integrato dipende dal suo utilizzo e dalla funzionalità che svolge.

in alcuni alimentaori esistono valorimi minimi di corrente vlaido, ovvero a cui possino ragigungere la tensione nominale. La maggior aprte ammette zero come valore minimo in quanto possono cmq ragigungere la tensione nominale anche senza avere connesso nessun carico.
Alcuni power supplies ay fail to regultaing the output voltage if the current is below a minumun, SWPS ricoadono in questa categoria.

Altri valori da prendere in considerazione per l'analisi della corrente in scuta sono i vlaori di massimo contiuboa, oivveor la corrent emasima che puòe essere. messa senza interruzni e quella di picco ovvero il valore massimo che può essere ragigunto per un ridotto periodo di tempo in modo ripetitivo.

Inoltre bisogna porre dei limiti sul valore massimo con cui la corrente può variare in unscita, in altre parole il massimo di:

\begin{equation}
	\frac{\unit{d}I_{out}}{\unit{d}t}
\end{equation}

questo pone un massimo alla variazione di corrente, nel caso in cui sia troppo grnade esse può causare alcuni prboemi come dei transienti di tensione troppo grandi e non desiderati.

Infine bisogna anche considerare i comprotamnti che il cirucito deve attuare nel caso in cui si verifichi un corto.
Più precimante in alcuni casi la corrente viene frazionata e limitata ad un certo vlaore current limiting oppure si entra nella fase di foldback ovvero si ha un'altra limitaizone della corrente, entrami i metodi sono attuatu al fine della protezione dei dispositvi alimentati ed in generale per un questione di sucurezza.

\subsubsection*{Power efficiency}


\subsubsection*{Protection}

\subsubsection*{Electromagnetic compatibility}

\subsubsection*{Ambient}

\subsubsection*{Mechanical}

\subsubsection*{Reliability and safety}


\subsubsection{Other specifications}









\end{document}








%\begin{figure}[h!]
%\begin{center}
%
%\begin{tikzpicture}[transform shape]
%	\draw (3, 16) to[american voltage source, l_={$V_s$}, i=$I_s$] (3, 14);
%	\draw (6, 17) to[american resistor, l={$R$}] (6, 15);
%	\draw (6, 15) to[cute inductor, l={$L$}] (6, 13);
%	\draw (3, 14) -- (3, 13) -- (6, 13);
%	\draw (3, 16) -| (3, 17) -- (6, 17);
%\end{tikzpicture}
%	
%\caption{Classical power supply}	
%\end{center}
%\end{figure}





%\begin{figure}[h!]
%\begin{center}
%
%\begin{tikzpicture}[transform shape, every node/.style={font=\bfseries\sffamily}]
%	\node[shape=rectangle, draw, line width=1pt, minimum width=1.965cm, minimum height=0.965cm] at (6, 8){} node[anchor=center, align=center, text width=1.577cm, inner sep=6pt] at (6, 8){Power\\processor\\};
%	\draw[line width=0.8pt, -latex] (8, 8) -- (9, 8);
%	\node[shape=circle, draw, line width=1pt, minimum width=0.715cm] at (9.375, 8){};
%	\node[shape=rectangle, draw, line width=1pt, minimum width=1.965cm, minimum height=0.59cm] at (6, 6.313){};
%	\draw[line width=0.8pt] (3, 8) -- (5, 8);
%	\node[circ] at (8, 8){};
%	\draw[line width=0.8pt, -latex] (8, 6.5) -- (7, 6.5);
%	\draw[line width=0.8pt, -latex] (6, 6.625) -- (6, 7.5);
%	\draw[line width=0.8pt] (8, 8) -| (8, 6.5);
%	\draw[line width=0.8pt, -latex] (8, 6.125) -- (7, 6.125);
%	\node[shape=rectangle, minimum width=3.215cm, minimum height=0.965cm] at (6, 6.313){} node[anchor=center, align=center, text width=2.925cm, inner sep=4.6pt] at (6, 6.313){Controller\\};
%	\node[shape=rectangle, minimum width=0.965cm, minimum height=0.965cm] at (9.375, 7.375){} node[anchor=center, align=center, text width=0.577cm, inner sep=6pt] at (9.375, 7.375){Load\\};
%	\node[shape=rectangle, minimum width=1.965cm, minimum height=0.965cm] at (9, 6.125){} node[anchor=center, align=center, text width=1.577cm, inner sep=6pt] at (9, 6.125){Reference};
%	\draw (3, 8) to[short, i^>=$i_i$] (5, 8);
%	\node[shape=rectangle, minimum width=0.965cm, minimum height=1.308cm] at (3.5, 8.25){} node[anchor=center, align=center, text width=0.577cm, inner sep=6pt] at (3.5, 8.25){$v_i$};
%	\draw[line width=0.8pt] (7, 8) -- (8, 8);
%	\draw (7, 8) to[short, i^>=$i_o$] (8, 8);
%	\node[shape=rectangle, minimum width=0.965cm, minimum height=1.762cm] at (8.125, 8.227){} node[anchor=center, align=center, text width=0.577cm, inner sep=6pt] at (8.125, 8.227){$v_o$};
%\end{tikzpicture}
%\caption{Block diagram of a power electronic system}
%\end{center}
%\end{figure}